\SourceLeaf{403}%
finden sich zahlreiche Flüchtigkeiten, die verbessert werden mussten; so fehlt im Besondern in der Formel 20 das erste Glied $\tfrac12\left(\frac1P\frac{dP}{du}-\frac1R\frac{dR}{du}\right)$. Vielleicht rührt das daher, dass \textsc{Gauss}, wie eine ausgestrichene Stelle zeigt, ursprünglich die Variabeln $t,\theta$ so gewählt hatte, dass
\[
\begin{aligned}
P&=aa+bb+cc=1\\
R&=\alpha\alpha+\beta\beta+\gamma\gamma=1
\end{aligned}
\]
ist, sodass also die Linien $t=\text{const.}$ und $\theta=\text{const.}$ geodätische Linien werden. In diesem Falle nimmt die Gleichung 21 die einfache Form an:
\[
\frac{df}{du}=\frac{d\chi}{du}+\frac{\theta't'}{ss}\left\{\cot\chi.Q't'-\cot(\omega-\chi)\left(\frac{d\omega}{du}h+Q''\theta'\right)\right\}.
\]
\begin{flushright}\textsc{Stäckel.}\end{flushright}
